\honest{Created Already In Motion}{Eliezer Yudkowsky}{2008}

Lewis Carroll wrote a short dialogue, ``What the Tortoise said to Achilles.'' Read it if you have not. The Tortoise takes two premises from Euclid, A and B, and a conclusion, Z. It accepts A and B but not the hypothetical that Z follows from them. Achilles asks it to accept that hypothetical as C. It does, and still will not accept Z, so Achilles asks it to accept D, and so on. Douglas Hofstadter's version shows the Tortoise turning each rule of inference ``into a mere string of the system.'' \nb{The quotations of Carroll and Hofstadter are accurate. Hofstadter's lines come from an exercise in \textsc{G\"odel, Escher, Bach}.}

My lesson: the Tortoise's mind needs the dynamic of adding Y to its beliefs when X and X$\to$Y are there. A rock lacks it, and you can add premises to a rock forever. A mind must be ``created already in motion.'' ``There is no computer program so persuasive that you can run it on a rock.'' \nb{Gilbert Ryle drew the same lesson from Carroll in 1946: knowing a rule of inference ``is not possessing a bit of extra information but being able to perform an intelligent operation.'' Yudkowsky had also made the point, without Carroll's name, in ``A Priori.''}

The same goes for action. A mind that believes a toddler is on the tracks, and that pulling toddlers off the tracks is ``fuzzle,'' does nothing unless it also implements a dynamic: when something is fuzzle, send it to the action system. A dynamic, I later added, is something that happens in a system, not data it stores; write one on paper and ``the paper will just lie there.'' A belief that sending fuzzle plans to action is fuzzle will not help, and neither will a proof that minds with the dynamic are more fuzzle, unless the mind already swaps in code it believes more fuzzle. \nb{This is Hume's view that ``reason alone can never be a motive to any action of the will.'' Nick Bostrom later used it to support the thesis that intelligence and goals can vary independently.} \nb{The post belongs to a sequence on metaethics, and ``fuzzle'' stands in for moral words. Some philosophers, the motivational internalists, hold that a sincere judgment that an act is right carries some motivation to do it, so a mind that does nothing about the toddler does not really judge rescue right. Others deny it. Either way, an AI must be built to act on its verdicts.}

``This is why you can't argue fuzzleness into a rock.''
